\documentclass[10pt,a4paper]{amsart}
\usepackage[margin=19mm]{geometry}
\usepackage{amsmath,amssymb,mathtools}
\usepackage[scr=rsfs,cal=euler]{mathalfa}
\usepackage{xfrac}
\usepackage[unicode,hidelinks,bookmarks=false]{hyperref}
\newcommand{\ie}{i.e.\ }
\newcommand{\cf}{cf.\ }
\newcommand{\resp}{respectively\ }
\newcommand{\Subsection}{\subsection}
\providecommand{\nobreakdash}{\nobreak\hbox{-}}
\setlength{\emergencystretch}{2em}
\multlinegap0pt
\usepackage{amssymb}
\newtheorem{thm}{Theorem}
\DeclareMathOperator{\cl}{cl}
\title{Finite Cardinalities of Mis\`ere quotients}
\author{Simon Rubinstein-Salzedo, Stephen Zhou}
\date{Source: arXiv:2603.18244v1 (CC BY 4.0)}
\begin{document}
\maketitle

\noindent\emph{Source and licence.} Finite Cardinalities of Mis\`ere quotients. Version arXiv:2603.18244v1 (CC BY 4.0), distributed by arXiv under the Creative Commons Attribution 4.0 International licence, http://creativecommons.org/licenses/by/4.0/, as recorded in the arXiv licence metadata for this exact version and shown on the abstract page https://arxiv.org/abs/2603.18244v1. Full licence notice: \texttt{licenses/arxiv-2603.18244-CC-BY-4.0.txt}.\par\smallskip

\noindent\textit{Editorial scope.} Counted unit: the article's thm thm (source0.tex lines 299--301). The article's complete original proof of it is delivered with it (source0.tex lines 302--324, one \texttt{proof} environment of 23 lines). No mathematics is written, repaired or re-derived: the statement and the proof are the article's own lines, in the article's own notation, and no retained line is substituted or re-flowed.  Retained uncounted as the source-local context the proof needs: lines 159--168; lines 171--184.  Nothing was omitted from any retained span, so the package carries no omission record.  No retained line contains a citation, so the wrapper carries no bibliography and no bibliography entry.  No retained line contains a figure environment, an image-inclusion command or drawing code, so no image is delivered and the package declares no asset dependency.  Editorial formatting only, in addition: an AMS \texttt{amsart} preamble that restates the article's own declarations and macros used by the retained lines, namely the environments \texttt{thm} on separate counters, together with the standard packages those lines need.  The retained environments print in this document as Theorem 1 in order of appearance, whereas the article prints the same items with the numbering of its own full document.  The complete native source of the article as distributed in the arXiv e-print for this exact version is bundled unchanged as \texttt{sources/196680/source0.tex}.
\begin{thm}
    \label{quotientsize}
    Let $H_i = \{x_{1,i},\dots,x_{2^{d_i},i}\}, 1 \leq i \leq m$ be a subspace of  $N_{<2^n}$ of dimension $d_i$, with $1 \in H_i$, and let $B_i = \{0,*\mid 0,*x_{1,i},\ldots,*x_{2^{d_i},i}\}$. Let $\mathscr{A} = \cl( *2^{n-1},B_1,\ldots,B_m ) $. If all subspaces $H_i$ are distinct and contain $1$, then $\left |\mathcal{Q} (\mathscr{A})\right | = \left( 2^n +\sum_{i=1}^m 2^{n-d_i} \right) +3 $ with the following equivalence classes:

i) the $2^n+2$ equivalence classes of $\mathscr{T}_n$,

ii) The $2^{n-d_i}$ equivalence classes $B_i + X$, where $\mathcal{G}(X) \equiv a \pmod {H_i}$ for a fixed $a \in N_{<2^n}$.

iii) A single equivalence class containing all sums in $\mathscr{A}$ with at least two blue mutant flowers.
\end{thm}

\begin{thm}
\label{quotientsize2}
    Let $\mathscr{A}$ be as in Theorem \ref{quotientsize}. Then $\left|\mathcal{Q} ( \mathscr{A} \cup \{2B_1 | \varnothing\})\right|  = \left |\mathcal{Q} (\mathscr{A})\right |+1$.

    The equivalence classes are as follows: 

    i) from Theorem \ref{quotientsize}

    ii) from Theorem \ref{quotientsize}

    iii) The equivalence class iii) from Theorem \ref{quotientsize}, expanded to include all sums in $\cl(\mathscr{A} \cup \{2B_1 \mid \varnothing\})$ containing at least one copy of $\{2B_1 \mid \varnothing\}$, excluding $n\{2B_1 \mid \varnothing\}$.

    iv) One equivalence class containing all $n\{2B_1 \mid \varnothing\}$ for $n>0$.
\end{thm}

\begin{thm}
    Let $N\neq  1,2,3,4,5,6,7,11$. Then there exists $\mathscr{A}$ such that $|\mathcal{Q}(\mathscr{A})| = N$.
\end{thm}

\begin{proof}

Recall that $N_{<2^n}$ is defined as the set of integers strictly less than $2^n$ viewed as a vector space over $\mathbb{F}_2$, where vector addition is $\oplus$. Notice that for $d\geq 1$, there exist $\binom{n-1}{d-1}$ subspaces of dimension $d$ that contain $1$. Now consider the binary representation of $N -3$, say \begin{equation}
    N-3 = 2^{a_0} + 2^{a_1} + \dots + 2^{a_{m}},
\end{equation}
where $a_0> a_0 > \dots > a_{m}$. Assume that $a_0 \geq 2$ and $m >0$. Set $n = a_1$ and $d_i = n - a_{i}$ for $i>0$. Since all the $a_i$'s are unique and decreasing, the $d_i$'s are unique and increasing. 

Assume that $a_m > 0$, so $d_i < n$ for all $i$.
Since $d_i \geq 1$ unconditionally and the $d_i$'s are unique, there exist distinct subspaces $H_1,H_2,\dots H_m$ of $N_{<2^n}$ such that $1 \in H_i$ and $H_i$ has dimension $d_i$ for all $i$. Let $H_i = \{x_{1,i},\dots,x_{2^{d_i},i}\}$ and define $B_i = \{0,*|0,*x_{1,i},\dots,*x_{2^{d_i},i}\}$ for $1 \leq i \leq m$. Let $\mathscr{A} = \cl( *2^{n-1},B_1,\dots,B_m ) $. Then by Theorem \ref{quotientsize}, \begin{equation}
    \left |\mathcal{Q} (\mathscr{A})\right | = \left( 2^n +\sum_{i=1}^m 2^{n-d_i} \right) +3 = \left( 2^n +\sum_{i=1}^m 2^{a_i} \right) +3 = N-3+3=N.
\end{equation}

If $a_m = 0$, so that $d_m = n$, the binary representation of $N-4$ is \begin{equation}
   N-4 = 2^{a_0} + 2^{a_1} + \dots + 2^{a_{m-1}}.
\end{equation}
We can proceed as before to get a mis\`ere quotient with cardinality $N-1$. But if we applied Theorem \ref{quotientsize2}, then we get a mis\`ere quotient with cardinality $N$ instead.

It remains to consider the case where $ m  =0$, so that $N = 2^{a_0} + 3$. Say that $a_0 > 3$, and let $n = a_0 -1$. Then $N_{<2^n}$ has at a subspace $H_1$ of dimension $1$ that contains $1$, and at least two subspaces $H_2,H_3$ of dimension $2$ that contain $1$.  Let  $B_i = \{0,*|0,*x_{1,i},\dots,*x_{2^{d_i},i}\}$ for $1 \leq i \leq 3$, and $\mathscr{A} = \cl( *2^{n-1},B_1,B_2,B_3 ) $. By Theorem \ref{quotientsize}, we see that \begin{equation}
    |\mathcal{Q}(\mathscr{A})| = \left(2^{n} + 2^{n-1} + 2^{n-1} \right)+3 = 2^{n+1}+3 = 2^{a_0}+3 =N.
\end{equation}

So unless $N \leq 2^2+2 = 6$, so that $a_0 < 2$, or $N = 2^{a_0} + 3$ for $a_0 \leq 3$, we have found $\mathscr{A}$ such that $|\mathcal{Q}(\mathscr{A})| = N$. These exceptions are exactly $N = 1,2,3,4,5,6,7,11$.
\end{proof}

\end{document}
